% === Diagrama 1: Árbol de problemas (rediseño legible) ===
% Diseñador-TikZ — Fase 5.   Tikz-Optimizer — pase de pulido visual.
% Contenido: §2.1 (párrafo 6) del Investigador.
% Paleta institucional: azulUNAL (#0066B3), grisLabIA (#666666), verdeGCPDS (#2E8B57).
% Las librerías TikZ (positioning, arrows.meta, shapes.geometric, calc, fit, backgrounds)
% ya están cargadas en main.tex; no se declaran aquí.
%
% Geometría (ancho útil ~16 cm):
%   Ramas  y Raíces: 4 cajas, text width=3.3 cm, centros x = -6, -2, 2, 6  (paso 4 cm)
%   Subproblemas   : 3 cajas, text width=3.6 cm, centros x = -6, 0, 6
%   Tronco         : 1 caja,  text width=12 cm, centrada en x = 0
%   Capas verticales (y): ramas 7.5, tronco 3.5, subproblemas -0.5, raíces -4.5
%   Gap vertical tronco→ramas ≈ 2.15 cm (cumple margen >1.5 cm para visibilidad de flechas).
%   Flechas causales ascendentes: verticales limpias para C1→SP1 y C4→SP3;
%   diagonales limpias para C2,C3→SP2 (forman "V" simétrica convergente en SP2.south).
%   Rellenos diferenciados para seguridad en escala de grises: verdeGCPDS!30 (ramas),
%   azulUNAL!22 (subproblemas), grisLabIA!22 (raíces).

\begin{figure}[t]
\centering

\begin{tikzpicture}[
  font=\footnotesize,
  >={Stealth[length=2.6mm, width=2.0mm]},
  line width=0.8pt,
  % --- Nodos ---
  rama/.style={rectangle, rounded corners=4pt, draw=verdeGCPDS, fill=verdeGCPDS!30,
               text width=3.3cm, align=center, inner sep=4pt, line width=0.9pt,
               minimum height=1.7cm, font=\footnotesize},
  tronco/.style={rectangle, rounded corners=6pt, draw=azulUNAL!70!black, fill=azulUNAL,
                 text=white, text width=12cm, align=center, inner sep=7pt, line width=1.1pt,
                 font=\small\bfseries, minimum height=2.0cm},
  subprob/.style={rectangle, rounded corners=4pt, draw=azulUNAL, fill=azulUNAL!22,
                  text width=3.6cm, align=center, inner sep=5pt, line width=0.9pt,
                  font=\small\bfseries, minimum height=2.0cm},
  raiz/.style={rectangle, rounded corners=4pt, draw=grisLabIA, fill=grisLabIA!22,
               text width=3.3cm, align=center, inner sep=4pt, line width=0.9pt,
               minimum height=1.8cm, font=\footnotesize},
  % --- Flechas ---
  flecha/.style={->, line width=0.9pt, draw=grisLabIA!75!black},
  flechacadena/.style={->, line width=0.9pt, draw=azulUNAL!75!black,
                       densely dashed, >={Stealth[length=2.6mm, width=2.0mm]}},
]

% --- Ramas (consecuencias) — nivel superior ---
\node[rama] (E1) at (-6, 7.5) {\textbf{E1.} Egresados con competencias superficiales y dependencia acrítica de la IA};
\node[rama] (E2) at (-2, 7.5) {\textbf{E2.} Desaprovechamiento del potencial de la IA agéntica para personalizar la formación};
\node[rama] (E3) at ( 2, 7.5) {\textbf{E3.} Baja transferencia de innovaciones en IA educativa hacia aulas y sector productivo};
\node[rama] (E4) at ( 6, 7.5) {\textbf{E4.} Desalineación entre perfiles de egreso y demandas de la Industria 4.0/5.0};

% --- Tronco (problema central) — nivel intermedio ---
\node[tronco] (TR) at (0, 3.5) {Limitada capacidad del sistema educativo en ingeniería para fortalecer el aprendizaje profundo —motivación intrínseca, argumentación técnica y resolución de problemas auténticos— en entornos mediados por IA, con tecnologías que no alcanzan madurez suficiente para su transferencia.};

% --- Subproblemas (puente causa → problema) ---
\node[subprob] (SP1) at (-6, -0.5) {SP1: Diagnóstico y modelado multidimensional del aprendiz};
\node[subprob] (SP2) at ( 0, -0.5) {SP2: Arquitectura agéntica autónoma para la personalización};
\node[subprob] (SP3) at ( 6, -0.5) {SP3: Validación empírica y transferencia TRL 6/7};

% --- Raíces (causas estructurales) — nivel inferior ---
\node[raiz] (C1) at (-6, -4.5) {\textbf{C1.} Predominio de pedagogías expositivas centradas en la transmisión de contenidos};
\node[raiz] (C2) at (-2, -4.5) {\textbf{C2.} Fragmentación del ecosistema de herramientas de IA sin orquestación pedagógica};
\node[raiz] (C3) at ( 2, -4.5) {\textbf{C3.} Escasa cultura de validación empírica y transferencia en \emph{edtech} para ingeniería};
\node[raiz] (C4) at ( 6, -4.5) {\textbf{C4.} Heterogeneidad en competencias digitales y matemáticas de los estudiantes en Colombia};

% --- Flechas: raíces → subproblemas (4 rutas limpias, sin solapamientos) ---
% Mapeo: C1→SP1, C2→SP2, C3→SP2, C4→SP3.  SP2 recibe 2 causas (C2,C3)
% que convergen en su bottom-center formando una "V" simétrica.
% C1→SP1: vertical directa
\draw[flecha] (C1.north) -- (SP1.south);
% C2→SP2: diagonal corta (C2 a la izquierda del centro de SP2)
\draw[flecha] (C2.north) -- (SP2.south);
% C3→SP2: diagonal corta (C3 a la derecha del centro de SP2)
\draw[flecha] (C3.north) -- (SP2.south);
% C4→SP3: vertical directa
\draw[flecha] (C4.north) -- (SP3.south);

% --- Flechas cadena lateral: SP1 → SP2 → SP3 (horizontal, punteada) ---
\draw[flechacadena] (SP1.east) -- (SP2.west);
\draw[flechacadena] (SP2.east) -- (SP3.west);

% --- Flechas: subproblemas → tronco (casi verticales) ---
\draw[flecha] (SP1.north) -- ($(TR.south west) + (0.5, 0)$);
\draw[flecha] (SP2.north) -- (TR.south);
\draw[flecha] (SP3.north) -- ($(TR.south east) + (-0.5, 0)$);

% --- Flechas: tronco → ramas (desde el borde superior del tronco) ---
\draw[flecha] ($(TR.north west) + (0.5, 0)$) -- (E1.south);
\draw[flecha] ($(TR.north west)!0.335!(TR.north east)$) -- (E2.south);
\draw[flecha] ($(TR.north west)!0.665!(TR.north east)$) -- (E3.south);
\draw[flecha] ($(TR.north east) + (-0.5, 0)$) -- (E4.south);

\end{tikzpicture}
\caption{Árbol de problemas del proyecto. Las raíces (causas estructurales C1--C4) alimentan los subproblemas SP1--SP3, los cuales confluyen en el tronco (problema central) cuyas ramas (consecuencias E1--E4) materializan el impacto negativo sobre el sistema educativo y el sector productivo. Las flechas indican relaciones causales ascendentes. Las flechas punteadas horizontales denotan secuencialidad entre subproblemas.}
\label{fig:arbol-problemas}
\end{figure}